\subsection{Functional reconstruction of the oriented moment}
\label{iii:sec:restitucion-funcional-catalan}

The differential relation~\eqref{iii:eq:catalan-vacuum} admits an
inverse preserving the complete collection of moments. Its domain is
the constitutive function produced by the cycle sectors, with the
orientation of~\eqref{iii:eq:uv}. The incidences $90=3\cdot30$ and
$120=4\cdot30$ determine that function before its evaluation in the channels
$s_\pm=q_\pm^{30}$. The condition at the origin corresponds to the
vanishing of the moment as its contractive argument tends to zero.
This condition also determines the integration constants.

\begin{theorem}[Integral reconstruction and uniqueness]
\label{iii:thm:restitucion-funcional-catalan}
Let $f(s)=(1-s^3)/(1-s^4)$, for $0<s<1$, be the response
of~\eqref{iii:eq:trace-det}, and let $\mathcal D=s\,d/ds$.
There is a unique function $F\in C^2((0,1))$ satisfying
\begin{equation}
 \mathcal D^2F(s)=\frac{s(1+s)}{1+s+s^2}f(s),
 \qquad \lim_{s\downarrow0}F(s)=0.
 \label{iii:eq:restitucion-funcional-problema}
\end{equation}
It is given by
\begin{align}
 F(s)&=\int_0^s\log\frac{s}{t}\,
          \frac{1+t}{1+t+t^2}f(t)\,dt
       =\int_0^s\frac{\log(s/t)}{1+t^2}\,dt
       =\mathcal C_2(s),
 \label{iii:eq:restitucion-funcional-integral}\\
 f(s)&=\frac{1+s+s^2}{s(1+s)}\mathcal D^2F(s).
 \label{iii:eq:restitucion-funcional-inversa}
\end{align}
Its continuous extension to $s=1$ recovers $F(1)=G_{\rm Cat}$.
\end{theorem}
\begin{proof}
The factorization $1+t+t^2+t^3=(1+t)(1+t^2)$ gives
\[
 \frac{1+t}{1+t+t^2}f(t)=\frac1{1+t^2}.
\]
The integrand of~\eqref{iii:eq:restitucion-funcional-integral} is
integrable at zero, since
\[
 0\leq F(s)\leq\int_0^s\log(s/t)\,dt=s.
\]
Hence $F(s)\to0$. Differentiation on the open interval,
first on a truncated interval and then by dominated convergence,
gives
\[
 \mathcal DF(s)=\int_0^s\frac{dt}{1+t^2},
 \qquad \mathcal D^2F(s)=\frac{s}{1+s^2}
     =\frac{s(1+s)}{1+s+s^2}f(s).
\]
In particular, $F$ is of class $C^2$ and satisfies the equation.
If $F_1,F_2$ are two solutions, their difference $H=F_1-F_2$
satisfies $\mathcal D^2H=0$. Setting $z=\log s$ gives
$d^2H(e^z)/dz^2=0$, so $H(s)=a\log s+b$.
The existence of a zero limit as $s\downarrow0$ first forces
$a=0$ and then $b=0$. The condition on $F$ therefore determines
the solution and also the limit $\mathcal DF(s)\to0$.

For $s<1$, the geometric expansion of $(1+t^2)^{-1}$ and
absolute integrability of its weighted terms allow termwise integration:
\[
 \int_0^s t^{2k}\log(s/t)\,dt
     =\frac{s^{2k+1}}{(2k+1)^2},\qquad
 F(s)=\sum_{k\geq0}\frac{(-1)^ks^{2k+1}}{(2k+1)^2}.
\]
This is precisely the collection~\eqref{iii:eq:c2}. Absolute and uniform
convergence of this last series on $[0,1]$ justifies passage
to the endpoint $s=1$. For the integral, this passage is likewise
justified by the integrable bound $\log(1/t)$, after extending
the integrand by zero for $t>s$. Finally, solving the differential
equation for $f$ gives~\eqref{iii:eq:restitucion-funcional-inversa},
with positive denominators on $(0,1)$.
\end{proof}

The limit of the moment and the limit of its second derivative are
taken in their respective convergence classes. In particular,
\[
 \lim_{s\uparrow1}\mathcal D^2\mathcal C_2(s)=\frac12
\]
is the radial Abel limit of the differentiated series. Evaluation
of $\mathcal C_2(1)$ directly uses its absolutely convergent
series. This distinction preserves the same depth operator in
the interior and in its boundary readout.

\subsection{Compatibility of reconstruction with channels and speed}
\label{iii:sec:restitucion-canales-velocidad}

The theorem reconstructs a function from another previously constructed
function. Cubic inversion~\eqref{iii:eq:cubic}, by contrast, acts
on the evaluations $r_\pm$: it recovers the arguments $s_\pm$
within that fixed collection. Composition of the two operations preserves
the quarter-turn orientation, the labels $P_\pm$, and
the transport incidences. Thus $G_{\rm Cat}=\mathcal C_2(1)$
is the endpoint readout of a functional object that also participates
in the interior responses $r_\pm=f(s_\pm)$.

For $\mathcal V_m=\mathcal V\otimes I_{9^m}$ and
$\tau_m=\operatorname{Tr}/(2\cdot9^m)$, the speed readout
retains exactly the normalization of
\eqref{iii:eq:velocidad-traza-normalizada}:
\begin{equation}
 \exp[-2\tau_m(\log\mathcal V_m)]
   =\frac1{r_+r_-}
   =(\det\mathcal V)^{-1}=\widehat c.
 \label{iii:eq:restitucion-velocidad-normalizada}
\end{equation}
Indeed, each eigenvalue $r_\pm$ is repeated $9^m$ times, so
$2\tau_m(\log\mathcal V_m)=\log r_++\log r_-$.
The identity uses the determinant of the two-sheet realization
and the normalized trace of its amplification. The previously fixed
dimensional sections then give
\[
 c_{\rm int}=c_{\rm vac}=\widehat c\,u_C,
 \qquad \ell_0=c_{\rm vac}t_0,
\]
with the same state, chart, and elementary duration as
in~\eqref{iii:eq:longitud-elemental-constitutiva}.
Functional reconstruction preserves this speed section when
composed with action and length; spectral multiplicity is
absorbed by the stated normalization.
